%%% LaTeX Template: Article/Thesis/etc. with colored headings and special fonts
%%%
%%% Source: http://www.howtotex.com/
% vim: set spell spelllang=es syntax=tex :

\documentclass[12pt]{article}
\usepackage[T1]{fontenc}
\usepackage{styles/apuntes-estilo}
\usepackage{styles/egyptian}
\usepackage{fancyhdr,lastpage}
\usepackage{hyperref}
\usepackage[inline]{enumitem}

\def\maketitle{

\makeatletter{
    \color{blue} \centering \huge \sc
    \textbf{
        Recuperatorio del primer examen parcial\\
        \small \vspace*{-8pt} \color{black}
        Tema: 3
        \vspace*{8pt}
    }\\
    \par
}
\makeatother

\makeatletter
% vim: set spell spelllang=es syntax=tex :
 {\centering \small 
    Introducción a la computación\\
    Departamento de Ingeniería de Computadoras \\
    Facultad de Informática - Universidad Nacional del Comahue \\
    \vspace{20pt} }
\makeatother

\vspace{-2.5cm}
\mbox{\hspace{-1cm}\includegraphics[width=3cm,height=3cm]{logos/uncoma.pdf}\hspace{12cm}
    \includegraphics[width=3cm,height=3cm]{logos/fai.pdf}}



}

% Custom headers and footers
\fancyhf{} % clear all header and footer fields
\fancypagestyle{plain}{\fancyhf{}}
\pagestyle{fancy}
\lhead{\footnotesize Recuperatorio del primer examen parcial}
\rhead{\footnotesize \thepage\ }

\def\ti#1#2{\texttt{#1} & #2 \\ }

\begin{document}

\thispagestyle{empty}
\maketitle
\setlength{\parindent}{1pt}

\textbf{Condiciones de la entrega:} La entrega del parcial se debe realizar en
la siguiente tarea de pedco:
\url{https://pedco.uncoma.edu.ar/mod/assign/view.php?id=350076}. Incluir las
hojas con los resultados y las resoluciones. De ser posible, nombre cada
fotografía con el formato \textbf{hoja\_\emph{<número de hoja>}}, o si se
entrega una foto por cada ejerció \textbf{ej\_\emph{<número de ejercicio>}}.
\textbf{Deben conservar las hojas originales}, ya que podemos pedirles que nos
envíen nuevas fotos.

\textbf{IMPORTANTE: \emph{Las resoluciones tienen que estar hechas en su puño
y letra.}} Está permitido tener impresas las tablas de la traza de ejecución,
pero deben estar completadas a mano.

\begin{enumerate}
\subsection*{Parte 1}

    \item ¿Qué número representa en base 10 el numeral $101$ si se lo
        interpreta en cada una de las siguientes bases?

        \begin{enumerate}

            \item Base 3.

            \item Base 5.

            \item Base 8.

        \end{enumerate}

    \item Utilizando una representación en \emph{punto fijo} y
        \emph{complemento a 2} utilizando 6 bits para la parte entera y 2 para
        la parte fraccionaria:

        \begin{enumerate}

            \item ¿Cuál es la representación de $25.57$?

            \item ¿Cuál es el error de representación (la diferencia entre el
                valor original y el almacenado)?

        \end{enumerate}

    \item Sean $A=-99_{10}$ y $B=86_{10}$:

        \begin{enumerate}

            \item Realice la suma $A+B$ en complemento a dos con 8 bits e
                \textbf{indique si se produce overflow o no}.

            \item Realice la resta $A-B$ en complemento a dos con 8 bits e
                \textbf{indique si se produce overflow o no}.

        \end{enumerate}

    \item ¿Cuál es la representación en hexadecimal utilizando el sistema
        \emph{Punto Flotante IEEE-754 de precisión simple (32 bits)} de
        $-25.5$ ?

\subsection*{Parte 2}

    \item ¿Cuál es la codificación en ASCII (representado en hexadecimal) del
        texto ``\textbf{ASCII usa 7b.}''? \emph{Respetar mayúsculas y
        minúsculas, e ignore las comillas}.

    \item Si se utiliza un esquema de compresión \emph{RLE(Run Lenght
        Encoding)} utilizando una codificación ``\emph{cantidad/color}'' de
        \textbf{2+2} (dos bits para la cantidad y dos bits para el color),
        donde la paleta de colores es 00 = blanco, 01 = amarillo, 10 = naranja
        y 11 = negro.

        \begin{enumerate}

            \item ¿Cuál es la imagen codificada en la siguiente secuencia de
                bytes representada en hexadecimal: \textbf{04 05 02 89 89 87
                46 4B E7}?

            \item ¿Cuántos bytes ocupa la imagen sin comprimir? \textbf{No
                olvide incluir la cabecera.}

        \end{enumerate}


\subsection*{Parte 3}

    \item Dado el siguiente programa:

        \begin{verbatim}
     0: START:  LD  I
     1:         ADD UNO
     2:         ST  I
     3:         ST  OUT
     4:         SUB MAX
     5:         JZ  FIN
     6:         JMP START
     7: FIN:    HLT
     8: I:      0
     9: MAX:    2
    10: UNO:    1
        \end{verbatim}

        \begin{enumerate}

            \item ¿Cuál es la salida del programa?

            \item ¿Cuál es el equivalente en código maquina binario de las
                instrucciones de las direcciones \textbf{0}, \textbf{5} y
                \textbf{6}?

        \end{enumerate}

    \item Dado el siguiente programa en binario:

        \begin{verbatim}
     0: 01001010
     1: 11100111
     2: 10101100
     3: 01101010
     4: 01001011
     5: 01111111
     6: 11011010
     7: 00100000
     8: 01111111
     9: 11011110
    10: 00000010
    11: 00000101
    12: 00000010
        \end{verbatim}

        ¿Cuál es la salida del programa?

    \item Sea un formato de instrucción que utiliza 5 bits para el código de
        operación y 11 para el operando que indica el número de la celda de
        memoria: ¿Cual es el máximo de instrucciones y celdas de memoria que
        se pueden acceder con este formato de instrucción?

\end{enumerate}

\subsection*{ \large\textbf{Anexo:} }

\textbf{Códigos de operación}

\begin{verbatim}
    010     LD
    011     ST
    100     ADD
    101     SUB
    110     JMP
    111     JZ
    001     HLT
    000     NOP
\end{verbatim}

\textbf{Tabla \emph{ASCII}:}

\begin{verbatim}
Dec Hex    Dec Hex    Dec Hex  Dec Hex  Dec Hex  Dec Hex   Dec Hex   Dec Hex
  0 00 NUL  16 10 DLE  32 20    48 30 0  64 40 @  80 50 P   96 60 `  112 70 p
  1 01 SOH  17 11 DC1  33 21 !  49 31 1  65 41 A  81 51 Q   97 61 a  113 71 q
  2 02 STX  18 12 DC2  34 22 "  50 32 2  66 42 B  82 52 R   98 62 b  114 72 r
  3 03 ETX  19 13 DC3  35 23 #  51 33 3  67 43 C  83 53 S   99 63 c  115 73 s
  4 04 EOT  20 14 DC4  36 24 $  52 34 4  68 44 D  84 54 T  100 64 d  116 74 t
  5 05 ENQ  21 15 NAK  37 25 %  53 35 5  69 45 E  85 55 U  101 65 e  117 75 u
  6 06 ACK  22 16 SYN  38 26 &  54 36 6  70 46 F  86 56 V  102 66 f  118 76 v
  7 07 BEL  23 17 ETB  39 27 '  55 37 7  71 47 G  87 57 W  103 67 g  119 77 w
  8 08 BS   24 18 CAN  40 28 (  56 38 8  72 48 H  88 58 X  104 68 h  120 78 x
  9 09 HT   25 19 EM   41 29 )  57 39 9  73 49 I  89 59 Y  105 69 i  121 79 y
 10 0A LF   26 1A SUB  42 2A *  58 3A :  74 4A J  90 5A Z  106 6A j  122 7A z
 11 0B VT   27 1B ESC  43 2B +  59 3B ;  75 4B K  91 5B [  107 6B k  123 7B {
 12 0C FF   28 1C FS   44 2C ,  60 3C <  76 4C L  92 5C \  108 6C l  124 7C |
 13 0D CR   29 1D GS   45 2D -  61 3D =  77 4D M  93 5D ]  109 6D m  125 7D }
 14 0E SO   30 1E RS   46 2E .  62 3E >  78 4E N  94 5E ^  110 6E n  126 7E ~
 15 0F SI   31 1F US   47 2F /  63 3F ?  79 4F O  95 5F _  111 6F o  127 7F DEL
\end{verbatim}

\end{document}
